\section{Results and discussion} \label{sec:results}
The results are organized in three parts. First, the synthetic experiment is used to assess the consistency of VAE training as well as to perform a sensitivity study to select an appropriate configuration (Sect.~\ref{sec:res:training}). The selected VAE configuration is then evaluated in terms of its ability to reproduce the statistical properties of the ensemble simulated by the dispersal model (Sect.~\ref{sec:res:evaluation}). For this purpose, a large test dataset comprising 2048 simulations was generated by FALL3D in the synthetic experiment, serving as a reference for evaluating the VAE-generated ensemble.  Finally, the Etna experiment demonstrates one possible application of this methodology to a real volcanic-cloud case: the latent-space assimilation of satellite observations (Sect.~\ref{sec:res:da}). The implications of these results and the main limitations of the proposed approach are discussed in Sect.~\ref{sec:discussion}.

\subsection{Training sensitivity analysis} \label{sec:res:training}
A sensitivity analysis was conducted to evaluate the performance of different neural network configurations trained with the synthetic experiment datasets. Multiple VAE configurations were trained by exploring a range of values for two key hyperparameters: the latent space dimension \latentDim and the factor $\beta$ in \eqref{eq:vae_loss}, while fixing the remaining hyperparameters with the values specified in Table~\ref{tab:hyperparameters}.
%
% VAE generation
%
For each trained VAE, an ensemble was generated by independently sampling latent vectors from the standard normal distribution and decoding them to recover the corresponding physical-space variables. In this section, each generated ensemble consisted of 2048 members.

%
% Multi-ensemble metrics
%
The performance was evaluated in terms of the RMSE using~\eqref{eq:rmse}, computed between the ensemble mean of the test dataset (taken as the reference) and the mean of the VAE-generated ensemble. To account for the stochastic variability of the VAE generation, 20 ensembles were generated for each combination of hyperparameters tested. The resulting RMSE values are shown as boxplots in Fig.~\ref{fig:boxplot}.
%
% Optimal hyperparameters
%
In all cases, the network was trained for 400 epochs, except for configurations with latent dimensions of $\latentDim=4$ and $\latentDim=2$, where early stopping was required to avoid overfitting. According to the trend suggested by Fig.~\ref{fig:boxplot}, the best performance is achieved for latent dimensions $\latentDim \gtrsim 6$ and $\beta\gtrsim 6$. Remarkably, performance degrades slightly for large latent space dimensions (Fig.~\ref{fig:boxplot}a). There are two possible explanations for this anomaly.
%
% Latent dimension discussion
%
First, large latent spaces often produce less stable training dynamics and worse generalization: enlarging latent spaces introduces more degrees of freedom than the data can realistically support. When the dimensionality grows beyond what is needed to represent the underlying factors, the model may spread information across latent variables in an inefficient way. This leads to poor compression, weaker regularization and a latent representation that becomes harder to structure during training, limiting the extraction of compact and meaningful representations. As a result, the decoder receives a less informative latent code, reducing the reconstruction quality. On the other hand, increasing the latent dimensionality can slow convergence: since the model must explore a much larger parameter space before reaching a stable configuration, it requires more training iterations and makes optimisation more difficult.

%
% Beta discussion
%
In addition, Fig.~\ref{fig:boxplot} shows the impact of the $\beta$ parameter, which controls the strength of the KL regularization in the loss function and, therefore, the balance between reconstruction fidelity and latent-space structure. Higher values force the model to impose a tighter, more organised latent space, at the expense of losing some fine-scale details in the reconstructions. In our case, $\beta \approx 6$ suggests that a strongly regularised latent representation is more beneficial for the task than capturing every small feature in the input data. This indicates that the model gains more from enforcing a well-structured latent space than from preserving high-resolution details during reconstruction.

\begin{figure}[htbp]
  \centering
  \includegraphics[width=0.4\textwidth]{f02}
  \caption{
  VAE sensitivity analysis for different combinations of two hyperparameters: the latent space dimension \latentDim (top) and the $\beta$ parameter (bottom), defined in \eqref{eq:L_vae}. 
  The performance is measured in terms of the Root-Mean-Square Error (RMSE), computed from the ensemble mean.
  }
  \label{fig:boxplot}
\end{figure}

%
% Latent space
%
This is further supported by the structure of the latent space shown in Fig.~\ref{fig:latent_space}. The nearly diagonal covariance matrix (Fig.~\ref{fig:latent_space}a) indicates weak correlation between the latent components, consistent with a well-regularized latent representation. Figure~\ref{fig:latent_space}b shows a two-dimensional visualization of the latent space based on the Uniform Manifold Approximation and Projection (UMAP) algorithm~\citep{mcinnes2020}, a technique for non-linear dimensionality reduction specifically designed to preserve the local structure of the data. The UMAP projection shows that the test dataset (circular markers) occupy the same region of latent space as the training data (plus-shaped markers), indicating that the learned representation generalizes to unseen data, rather than only organizing the training samples.
%
In addition, it is also worth showing how the encoding can discriminate between qualitatively different model solutions (ensemble members) by grouping them in different regions of the latent space. This is illustrated in Fig.~\ref{fig:latent_space}b by means of coloured-scale dots, where members with eruption column heights higher and lower than the reference (ensemble central member) are clearly concentrated in disjoint (\ie non-overlapping) regions of latent space. In other words, the VAE encoding implicitly produces a latent space with clustering-like structure.
%
% Conclusion
%
These results indicate that the network was trained correctly and behaves as expected. Based on this sensitivity analysis, a latent dimension of $\latentDim=16$ and $\beta=6$ were adopted for the VAE configuration used in the next section.

\begin{figure}[htbp]
  \centering
  \includegraphics[width=0.4\textwidth]{f03}
  \caption{%
    Structure of the latent space. %
    The independence between the distributions of the latent vector components can be confirmed from the covariance matrix (a). 
    The test dataset is encoded in the same region of the latent space covered by the training dataset, as shown in the UMAP two-dimensional visualization of the latent space (b), confirming that the learned representation generalizes well to unseen data.
    }
  \label{fig:latent_space}
\end{figure}

\subsection{Evaluation of the VAE-generated ensemble} \label{sec:res:evaluation}
Once the VAE neural network has been trained with numerical simulations, new samples can be generated at very low computational cost. In this section, we explore whether these samples can reproduce the statistical characteristics of an independent and larger ensemble of numerical simulations (test dataset from the synthetic experiment), even when the new samples have not been obtained by explicitly solving the underlying physical equations.
%
% VAE-generated ensemble in the physical space
%
For illustrative purposes, Fig.~\ref{fig:vae_ensemble} shows a few VAE-generated samples of two-dimensional ash column mass. The variability in the ensemble is a consequence of the different eruptive source parameters (\eg, eruptive column height) and meteorological conditions used to simulate the training dataset. For example, samples in Fig.~\ref{fig:vae_ensemble}c,d correspond to cases with high eruptive column heights and strong emission, whereas other samples are associated with low column heights and weak emission (\eg Fig.~\ref{fig:vae_ensemble}h). In most cases, the volcanic plume is advected eastwards from the vent (16.76\degree~W, 28.27\degree~N). Additionally, some cases are characterised by a secondary plume with a SE- or even W-component in the transport pattern, suggesting the presence of strong wind shear that introduces variability into the ensemble and leads to qualitatively different solutions.

\begin{figure*}[htbp]
  \centering
  \includegraphics[width=\textwidth]{f04}
  \caption{%
    The neural network was trained to generate samples of two-dimensional ash column mass, \ie the total mass of volcanic ash in a vertical column.
    In this example, a VAE-generated ensemble composed of 12 samples is shown.
    }
  \label{fig:vae_ensemble}
\end{figure*}

\subsubsection{Ensemble mean and spread}
The VAE ability to reproduce sample statistics, such as the ensemble mean and spread (standard deviation), is examined in this section and a few metrics are reported. Figure~\ref{fig:mean} compares the ensemble means of the training ensemble (Fig.~\ref{fig:mean}a) and VAE-generated ensemble comprising 5000 members (Fig.~\ref{fig:mean}b). Ideally, VAE should be able to closely approximate the test dataset (target) to achieve good performances. Deviations from the target are quantified by bias maps. In relative terms, the bias is small for both the training dataset (Fig.~\ref{fig:mean}c) and VAE (Fig.~\ref{fig:mean}d), although the training dataset provides a slightly better estimate of the ensemble mean in general.

\begin{figure*}[htbp]
  \centering
  \includegraphics[width=\textwidth]{f05}
  \caption{%
    Ensemble mean maps obtained from the training dataset (a) and a VAE-generated ensemble with 5000 samples (b). 
    The bias maps for the training (c) and VAE (d) ensemble means are computed by subtracting the target (test dataset) corresponding results.
    }
  \label{fig:mean}
\end{figure*}

A summary of the metrics is shown in Table~\ref{tab:metrics}. Metrics are computed for multiple VAE-generated ensembles composed of 5000 samples in order to account for the variability introduced by stochastic sampling. The average over 20 ensembles is reported along with the standard deviation for the uncertainty. The target test dataset is taken as a reference in order to evaluate each of the metrics described in Sect.~\ref{sec:metrics}. The VAE errors measured in terms of the RMSE are twice as high as those evaluated from the training dataset. In relative terms, this represents an error of 1.78\% (training) and $3.4\%$ (VAE) according to the MAPE evaluated from the ensemble mean. In order to avoid divisions by zero and numerical instabilities, the MAPE was actually computed considering only the region where the test dataset is above 5~\unit{g\,m^{-2}} (\ie the range shown in Figs.~\ref{fig:mean}a,b). We can conclude that VAE is properly capturing the first (mean) and second moments (variance) of the distribution, although the training estimations for the ensemble mean and spread are still slightly better.
%
The next section demonstrates that, when larger ensembles are generated, the VAE tends to perform better on statistics more sensitive to sampling errors.

\begin{table}[t]
  \caption{
    A summary of some performance metrics evaluated using the test dataset as the reference target.
    The VAE metrics were averaged over 20 ensembles generated to account for stochastic variability.
    The uncertainty range is given by the standard deviation.
  }
  \label{tab:metrics}
  \begin{tabular}{lccc}
    \tophline
    Metric    & Variable  & Training dataset  & VAE (5000 samples) \\ \middlehline
    RMSE      & Mean      & 0.108             & $0.20  \pm 0.03$   \\
    MBE       & Mean      & 0.001             & $-0.04 \pm 0.02$   \\
    MAPE (\%) & Mean      & 1.784             & $3.4   \pm 0.6$    \\ \middlehline
    RMSE      & Spread    & 0.255             & $0.56  \pm 0.04$   \\
    MBE       & Spread    & -0.004            & $-0.17 \pm 0.01$   \\
    \bottomhline
  \end{tabular}
  \belowtable{ 
    Values are reported in \unit{g~m^{-2}}, except MAPE, which is given as a percentage (\%).
    }
\end{table}

\subsubsection{Probability distributions} \label{sec:res:pd}
An important aspect in hazard and impact assessment is to evaluate exceedance probabilities, \ie the probability that a specific hazardous threshold is exceeded within a given time frame. In the context of ensemble modelling, probability distribution estimations can suffer from sampling errors, particularly when small ensembles are considered. Figure~\ref{fig:exceedance} presents the exceedance probabilities for three different thresholds, demonstrating that the VAE is capable of generating accurate probability distributions as well as realistic samples.
%
In order to examine whether the sampled ensemble provides a good approximation to the target distribution, a few histograms are shown in Fig.~\ref{fig:histograms} at four selected locations, comparing results for the training dataset (left column) and a VAE ensemble with 5000 samples (right column), considering 128 bins between the minimum and maximum values for each location. As expected, the distribution estimated from the training dataset with a small ensemble size ($\ensSize=256$) results in a very noisy distribution, while the VAE distribution provides a good approximation to the test dataset distribution, even for multimodal distributions (\eg, Fig.~\ref{fig:histograms}a).

\begin{figure*}[htbp]
  \centering
  \includegraphics[width=\textwidth]{f06}
  \caption{%
    Exceedance probabilities for three ash column mass thresholds (8, 16 and 20~$g~m^{-2}$ from top to bottom). 
    The probability contours correspond to $2\%$, $25\%$, $50\%$, $75\%$ and $98\%$.
    Results for the training dataset (left column) and a VAE-generated ensemble with 5000 samples (central column) are compared to the test dataset target (right column).
    }
  \label{fig:exceedance}
\end{figure*}

\begin{figure*}[ht]
  \centering
  \includegraphics[width=0.8\textwidth]{f07}
    \caption{
      Histograms for the training dataset (256 samples, left column) and a 
      VAE-generated ensemble (5000 samples, right column) at different locations. 
      The results are compared against the test dataset (2048 samples), 
      assumed to be the reference target (yellow shaded area). 
      For illustrative purposes, four sampling locations were selected: 
      P1 at 15.5\degree\,W, 26.3\degree\,N \textbf{(a, b)}; 
      P2 at 16.8\degree\,W, 25.9\degree\,N \textbf{(c, d)}; 
      P3 at 15.0\degree\,W, 27.3\degree\,N \textbf{(e, f)}; and 
      P4 at 15.5\degree\,W, 27.4\degree\,N \textbf{(g, h)}.
    }
  \label{fig:histograms}
\end{figure*}

Naturally, the VAE capacity to approximate the distribution depends on the number of samples generated. In order to assess the quality of the approximation for different amounts of samples, the JS divergence was computed from discrete probability distributions using~\eqref{eq:js}. Specifically, $JSD$ was obtained from histograms with a fixed range (128 bins between 0 and 40 $g~m^{-2}$) for each grid cell. As a measure of the worst-case conditions, the 99th percentile of the grid‐cell values is reported in Fig.~\ref{fig:js_divergence}, taken as a robust estimation of the maximum $JSD$ over the domain. The results for multiple VAE-generated ensemble sizes are compared with the training dataset. As expected, VAE improves the training distribution only for large ensemble sizes ($\ensSize\gtrsim 256$), reaching an asymptotic value of approximately $JSD \sim 0.042$, which represents a 2.6x improvement over the corresponding training metric ($JSD=0.11$).

\begin{figure}[htbp]
  \centering
  \includegraphics[width=0.45\textwidth]{f08}
  \caption{%
    The Jensen-Shannon (JS) divergence was computed for each grid cell 
    considering different sizes of the VAE-generated ensemble.
    The 99th percentile is reported as a worst-case measure of performance.
    Low values indicate a good approximation to the actual probability distribution. 
    The result for the training dataset is also indicated for reference.
    }
  \label{fig:js_divergence}
\end{figure}

\subsubsection{Isotropic power spectrum} \label{sec:psd}
%
% PSD, general description
%
The isotropic power spectrum, \ie the radially averaged two-dimensional Power Spectral Density (PSD), was computed using \eqref{eq:psd}. Figure~\ref{fig:psd} presents the results for the test dataset and VAE, averaged over multiple ensemble realizations (each ensemble with $\ensSize=5000$ samples). It is worth noting that, instead of the classical energy cascade spectrum proportional to $k^{-5/3}$, the test dataset exhibits an exponential decay (solid red line), likely due to the diffusion processes introduced by the dispersal model FALL3D. The generative model, with a strongly regularized latent space, reproduces well the large scales but fails at smaller scales, as shown in Fig.~\ref{fig:psd} (green triangles).
%
% High frequency noise
%
The isotropic power spectrum deviates from the exponential decay for wavenumbers above $\sim 0.3~\Delta x^{-1}$, approximately, where $\Delta x$ is the grid size. In other words, the VAE is unable to reproduce variability at spatial scales below $\sim 3~\Delta x$. The flat high-wavenumber tail in the 2D power spectrum seems to be related to the convolutional kernel size used in this work (\ie, $3\times 3$) and the limitations of the decoder upsampling. Our decoder architecture is based on transposed convolution layers, and it is widely known that transposed convolutions usually lead to chequerboard artifacts. Actually, this effect can also be observed from the bias map of the ensemble mean (Fig.~\ref{fig:mean}d), where a noisy pattern with relatively low intensity and high-frequency structure can be identified.
%
% Kernel size
%
A new neural network was trained using a larger kernel size ($5\times 5$) in order to investigate whether it yields to a more effective low-pass filtering (blue crosses in Fig.~\ref{fig:psd}). Although the background noise was reduced slightly, the results remained almost identical. Future work should investigate more advanced neural network architectures to better address this limitation.
%
% High beta
%
Finally, this behaviour is consistent with the disentangled latent representation enforced by a high value of the $\beta$ hyperparameter (see Fig.~\ref{fig:boxplot}). In fact, the VAE learns a smooth, structured encoding that preserves the dominant generative factors while filtering out small-scale variations that would otherwise help reduce the reconstruction error. The effects of this strongly regularized latent space are also visible in the spectral analysis. While the VAE successfully reproduces the low-frequency components associated with large-scale structures, it systematically underestimates the high-frequency content responsible for fine-scale details. 

\begin{figure}[htbp]
  \centering
  \includegraphics[width=0.45\textwidth]{f09}
  \caption{%
    The isotropic power spectrum, \ie radially averaged 
    two-dimensional Power Spectral Density (PSD), 
    reveals inconsistencies in the spatial structure of 
    the VAE-generated samples. 
    The generative model reproduces large scales but fails 
    at scales smaller since the VAE is introducing 
    artificial noise for large k. Since the wavenumber 
    is expressed in terms of the grid size, $\Delta x$, 
    we can conclude that the VAE is unable to reproduce 
    the variability over spatial scales below 
    $\sim 3~\Delta x$, approximately. 
    }
  \label{fig:psd}
\end{figure}

\subsection{Latent-space assimilation of satellite observations} \label{sec:res:da}
This section shows the results of the latent-space data assimilation for the Etna experiment (see Sect.~\ref{sec:etna-experiment}). Figure~\ref{fig:etna2018}a shows the ensemble mean of the model prior ($\ensSize=128$), \ie the ensemble of two-dimensional ash column mass at 18:00~UTC simulated with the dispersal model FALL3D. The corresponding satellite retrievals from SEVIRI are shown in Fig.~\ref{fig:etna2018}b. The relatively complex spatial structure of the observed plume results from a pronounced clockwise rotation (veering) of its trajectory during the initial phase of the eruption. This change in flow direction is not captured by the driving meteorological dataset, which leads FALL3D to simulate ash transport under a largely uniform directional flow. Consequently, this case represents a challenging test for the data assimilation methodologies.

\begin{figure}[htbp]
  \centering
  \includegraphics[width=0.9\textwidth]{f10}
  \caption{%
    Spatial distribution of ash column mass at 18:00~UTC for the Etna 24 December 2018 case for 
    \textbf{(a)} the ensemble mean of the model prior, 
    \textbf{(b)} the SEVIRI satellite retrievals, 
    \textbf{(c)} the ETKF posterior ensemble mean, and 
    \textbf{(d)} the PF posterior ensemble mean. 
    The red box shows the FALL3D model domain extent.
    }
  \label{fig:etna2018}
\end{figure}

\subsubsection{VAE prior}
The simulated model prior is used to train a new VAE, which is subsequently used to generate an augmented ensemble (VAE prior) by sampling from the latent space. The same procedure and neural network architecture used for the synthetic experiments are employed for the Etna experiment, with a few exceptions:
(i) to accommodate the increased size of the input data, the latent dimension was increased to $\latentDim=32$ and a new convolutional layer was added to both the encoder and decoder;
(ii) we found that the checkerboard pattern can be reduced by replacing transposed convolutional upsampling with bilinear interpolation; and 
(iii) rather than sampling the latent variables from a standard normal distribution, the latent vectors are sampled from a multivariate normal distribution,
\begin{equation}
\mathbf{z}_i \sim \mathcal{N}(\boldsymbol{\mu}, \boldsymbol{\Sigma}),
\end{equation}
where $\boldsymbol{\mu}$ and $\boldsymbol{\Sigma}$ are estimated from the latent representations of the training ensemble. The VAE prior is then generated by sampling $\ensSize=4096$ latent vectors $\mathbf{z}_i$ from this distribution and then decoded to obtain the VAE prior in the original ash column mass space.

The agreement with the satellite observations is evaluated using the metrics defined in Sect.~\ref{sec:metrics}. Table~\ref{tab:rmse_metrics} compares the model prior and VAE prior against the satellite observations. The metrics are computed for every ensemble member and the ensemble average is reported. Both model and VAE ensembles show very similar performance across all metrics, with the VAE prior yielding slightly better RMSE, RMSLE and Pearson correlation. These results indicate that the VAE-generated ensemble preserves the main characteristics of the model prior relevant to the comparison with observations. The results are also consistent with the previous analysis, showing that the VAE can generate samples that are statistically comparable to those produced by the physics-based FALL3D model.
%
% Concluding remarks
%
These results provide a baseline for assessing the impact of data assimilation in the next section, while also demonstrating that the VAE-generated ensemble provides a suitable prior for subsequent data assimilation experiments.

\begin{table}[htbp]
    \centering
    \caption{Metrics used to quantify agreement with observations for the Etna data assimilation experiment.}
    \label{tab:rmse_metrics}
    \begin{tabular}{lcccc}
        \tophline
         & Prior   & Prior & Posterior & Posterior \\
         & (model) & (VAE) & (ETKF)    & (PF) \\ \middlehline
        RMSE & 180.317 & 173.562 & 145.696 & 126.903 \\
        RMSLE & 1.561 & 1.491 & 1.228 & 1.210 \\
        Pearson correlation & 0.315 & 0.339 & 0.700 & 0.667 \\
        POD & 0.278 & 0.274 & 0.381 & 0.378 \\
        Jaccard index & 0.229 & 0.227 & 0.325 & 0.333 \\
        \bottomhline
    \end{tabular}
\end{table}

\subsubsection{Data assimilation experiments}
%
% Spatial results
%
The analysis ensemble means obtained with the ETKF (Fig.~\ref{fig:etna2018}c) and PF (Fig.~\ref{fig:etna2018}d) show that both the intensity and location of the maximum are corrected relative to the prior. However, the analysis ensembles retain the lack of representation of the veering and wind shear present in the prior, resulting in an almost unidirectional and less extended plume. This illustrates an important limitation of the assimilation: the analysis can preferentially select and combine ensemble members that are more consistent with the observations, but cannot recover spatial structures that are not represented in the prior ensemble.
%
% Quantitative improvement
%
The impact of the assimilation can be assessed in terms of the evaluation metrics reported in Table~\ref{tab:rmse_metrics}. These diagnostics are not intended as an independent validation, but rather to characterize the adjustment of the prior ensemble towards the observational constraints. The agreement with the observations improves for all metrics following assimilation. The RMSE is reduced by 16\% (ETKF) and 27\% (PF) relative to the VAE prior, while the RMSLE decreases by 18\% (ETKF) and 19\% (PF). The Pearson correlation approximately doubles for both methods. A larger relative improvement is obtained for the categorical metrics: the ETKF increases the POD and Jaccard index by 39\% and 43\%, respectively, while the corresponding increases for the PF are 38\% and 47\%. Overall, both assimilation approaches consistently adjust the ensemble towards the observations, with comparable performance across the different metrics considered.

%
% Ensemble spread
%
The impact of the assimilation can also be assessed in terms of the ensemble spread. The ensemble spread quantifies the dispersion of the ensemble members around the mean and is computed as the ensemble standard deviation at each element, averaged over all elements. The resulting spreads for the prior and posterior ensembles are reported in Table~\ref{tab:spread_metrics}. This metric is evaluated both in physical space and in latent space.
%
% PF collapse
%
The spread reduction in latent space is 10\% (ETKF) and 42\% (PF). The substantially lower PF spread results from particle-filter degeneracy, with nearly all posterior probability concentrated on only two particles. Kernel perturbation is subsequently used to increase ensemble diversity, with a regularization factor of $h=0.2$ (see Sect.~\ref{sec:pf}). Interestingly, the spread reduction in physical space is more pronounced for the ETKF (77\%) than for the PF (70\%), despite the substantially larger reduction observed for the PF in latent space.

\begin{table}[htbp]
    \centering
    \caption{Ensemble spread of the prior and posterior ensembles in latent and physical space}
    \label{tab:spread_metrics}
    \begin{tabular}{lrr}
        \tophline
        ~              & Latent-space spread & Physical-space spread \\ \middlehline
        VAE Prior      & 0.999 & 33.974 \\
        ETKF Posterior & 0.899 & 7.804 \\
        PF Posterior   & 0.582 & 10.351 \\
        \bottomhline
    \end{tabular}
\end{table}

%
% Latent-space distributions
%
The latent-space distributions of the prior and posterior ensembles are shown in Fig.~\ref{fig:latent-distribution} for selected pairs of latent dimensions, together with their corresponding histograms. In general, the prior distribution is well approximated by the standard normal distribution, as expected from the VAE construction. The ETKF produces a Gaussian-like posterior, whereas the PF analysis is clearly bimodal, reflecting the concentration of the posterior weights on two particles during the initial filter update. The subsequent kernel perturbation broadens these two modes, restoring some variability.

\begin{figure}
    \centering
    \includegraphics[width=0.42\linewidth]{f11.png}
    \caption{
        Latent-space distributions of the prior and posterior ensembles from the data assimilation experiments. Three selected pairs of latent dimensions are shown for illustration. The ensembles, each comprising 4096 members, are: VAE prior (white), ETKF posterior (orange), and PF posterior (green).
        }
    \label{fig:latent-distribution}
\end{figure}


\subsection{Discussion} \label{sec:discussion}
%
% General remarks
%
In this work, the VAE is viewed as a complement to the physics-based model. Making use of an ensemble of a few hundred simulations, the VAE learns a compact, low-dimensional representation of the complex, high-dimensional variability captured by the physics-based simulations. This representation provides an alternative space in which the structure of the ensemble can be analysed and manipulated. 

%
% Ensemble augmentation
%
Sampling in the latent space to generate an augmented ensemble is the most straightforward application of this methodology. Although it does not introduce new physical knowledge, the VAE can generate a much larger number of realizations at very low computational cost during inference, while retaining variability that is statistically consistent with that captured by the physics-based simulations.
%
% Computational cost
%
It is interesting to consider whether this methodology can be computationally efficient, accounting for both training and inference, when generating ensembles larger than the training ensemble. However, the answer depends on the available computational resources, the cost of the physics-based model, the total simulation time, and the ensemble size, among other factors. As an illustrative example, let's consider the Etna experiment, which is the most computationally demanding case presented here. The 128 FALL3D simulations were run on the accelerated partition of MareNostrum5 using 128 NVIDIA Hopper GPUs, requiring approximately 1 GPU-hour in total. In contrast, VAE training was performed on a laptop using a single RTX 5000 Ada Generation NVIDIA GPU and required only $5.83 \times 10^{-3}$ GPU-hours. In this experiment, the forecast period was limited to 8~h; extending the forecast period would further increase the relative cost of the numerical simulations. Most importantly, the approach substantially reduces the computational resources required to generate large ensembles; for example, the simulations for the synthetic experiment on the LUMI supercomputer required at least 2,048 GPUs. Of course, these examples do not allow us to draw general conclusions about the relative efficiency of augmenting an ensemble with a VAE versus explicitly simulating larger ensembles, as the computational cost depends on the type of model and the available computational resources. However, computational efficiency is arguably not the main advantage of the proposed approach.

%
% Sampling errors
%
This methodology presented in this work also offers potential advantages beyond ensemble augmentation. Increasing the ensemble size can reduce sampling errors and improve the estimation of statistical quantities that are particularly sensitive to small ensemble sizes, such as exceedance probabilities and other measures of uncertainty. Such estimates are particularly relevant for probabilistic hazard assessment and decision-making in crisis management. For example, as shown in Sect.~\ref{sec:res:pd}, the VAE provides a better approximation to the reference distribution than the training ensemble in terms of the Jensen--Shannon (JS) divergence when the ensemble size exceeds the number of numerical simulations used for training.

%
% Sample generation
%
Additionally, due to the generative nature of the VAE, instead of simply duplicating or resampling existing numerical simulations, as in a bootstrap approach, new realizations absent from the original dataset can be generated. This could potentially lead to a richer statistical representation of the variability captured by the physics-based model, particularly if some of the generated realizations correspond to regions of the distribution that are poorly represented in the training ensemble. This interpretation is consistent with the results presented in Sect.~\ref{sec:da}, where the VAE-generated ensemble shows better agreement with the observations than the model prior in terms of RMSE, RMSLE, and Pearson correlation (Table~\ref{tab:rmse_metrics}).

%
% Operations in the latent space
%
Most importantly, the learned representation enables different operations to be performed in a substantially lower-dimensional space, including data assimilation, clustering, perturbation, conditional sampling, and other statistical techniques. These capabilities provide an excellent opportunity for further research. In this study, we explore two latent-space data assimilation approaches that, in the formulation adopted here, would not be directly applicable in the original physical space.
%
% Data assimilation
%
The assimilation experiments demonstrate the feasibility of performing data assimilation directly in the VAE latent space. Both the ETKF and PF adjust the augmented prior towards the satellite observations, improving the agreement across all considered evaluation metrics and reducing the ensemble spread without an excessive loss of ensemble variability. Although the two assimilation approaches produce markedly different posterior distributions in latent space, their resulting posterior ensembles in physical space are relatively similar.
%
% Particle filter collapse
%
The particle-collapse problem of the PF was not fully resolved, as the posterior distribution remained concentrated around two particles, resulting in a bimodal distribution after the perturbation introduced by the kernel regularization. Applying such a naive perturbation directly in the physical space for a nonlinear model would result in states that are inconsistent with the model dynamics. Instead, perturbing in latent space takes advantage of the learned mapping, allowing perturbed particles to be mapped back to physical space while retaining consistency with the physics-based simulations. This clearly illustrates one of the advantages of performing operations directly in the latent space.

Overall, the results presented in Sect.~\ref{sec:res:da} demonstrate that latent-space data assimilation can effectively incorporate satellite observations to correct a large VAE-generated ensemble of two-dimensional volcanic ash column mass. Future research is needed to assess whether the methodology can be extended to the full three-dimensional concentration field and multiple particle classes, which will require adapting the neural network architecture to the increased dimensionality of the problem. This extension would be an important step for enabling multiple assimilation cycles and, ultimately, supporting operational forecasting.